Este anexo reúne las siglas y los símbolos que se usan en el informe. Los términos del problema se definen en la sección~\ref{sec:terminos}.

\begin{table}[H]
\centering
\caption{Siglas.}
\label{tab:siglas}
\small
\begin{tabularx}{\textwidth}{@{}l X@{}}
\toprule
Sigla & Significado \\
\midrule
RED & Red Metropolitana de Movilidad de Santiago \\
GTFS & \emph{General Transit Feed Specification}, formato estándar de horarios de transporte público \\
E-VSP & Programación de vehículos eléctricos (\emph{electric vehicle scheduling problem}) \\
MILP & Programación lineal entera mixta (\emph{mixed-integer linear programming}) \\
KPI & Indicador de desempeño (\emph{key performance indicator}) \\
E1 a E5 & Electroterminales Vespucio Norte, El Conquistador, Los Espinos, La Reina y Santa Rosa \\
T01 a T04 & Terminales comunes, sin cargadores \\
L, S, D & Día laboral, sábado y domingo \\
\bottomrule
\end{tabularx}
\end{table}

\begin{table}[H]
\centering
\caption{Símbolos del modelo del capítulo~\ref{ch:modelo}.}
\label{tab:simbolos}
\small
\begin{tabularx}{\textwidth}{@{}l X@{}}
\toprule
Símbolo & Significado \\
\midrule
\(N\) & Corridas del tipo de día \\
\(T,\ E,\ F\) & Terminales comunes, electroterminales y lugares de término (\(F=T\cup E\)) \\
\(s_i,\ f_i,\ L_i,\ \sigma(i)\) & Inicio, término, distancia de servicio y sector de la corrida \(i\) \\
\(A\) & Pares de corridas que un bus puede encadenar directo \\
\(K_{ij},\ R\) & Opciones de carga entre \(i\) y \(j\), y conjunto de todas ellas \\
\(e(k),\ a_k,\ h_k,\ g_k,\ p_k\) & Electroterminal, km de ida y de vuelta, energía máxima y precio medio de la opción \(k\) \\
\(d_{ij},\ d^{0}_i,\ d^{1}_{i\ell}\) & Km del traslado entre corridas, desde el terminal de salida y hasta el lugar de término \\
\(B,\ m,\ r\) & Batería, reserva y consumo por km \\
\(c^{B},\ c^{D},\ c^{S}\) & Costo por bus, por km y por sesión \\
\(P_e,\ C_e\) & Plazas y conexiones del electroterminal \(e\) \\
\(\Omega,\ \alpha_{k\omega}\) & Bloques de cinco minutos e indicador de que la opción \(k\) carga en el bloque \(\omega\) \\
\(u_i,\ z_{i\ell}\) & La corrida \(i\) inicia una jornada, o la termina en \(\ell\) \\
\(x_{ij},\ y_{ijk},\ q_{ijk}\) & Encadenamiento directo, encadenamiento con carga y kWh cargados \\
\(e_i\) & Energía del bus al terminar la corrida \(i\) \\
\bottomrule
\end{tabularx}
\end{table}
